\section{Conclusion}

We introduced mmIR, a differentiable inverse-rendering framework that fits scene and sensor parameters directly from raw FMCW ADC. By combining end-to-end automatic differentiation through phase-coherent MIMO path accumulation with ITU physics-based per-vertex materials, mmIR optimizes materials, normals, and beam patterns from real radar captures. On seven outdoor scenes, mmIR achieves 0.63 mean Pearson correlation on range--azimuth maps, substantially outperforming Sionna-RT (0.32). After fitting, the optimized scene produces sensor-realistic ADC from dense virtual apertures for 3D occupancy detection and scene reconstruction, validated against LiDAR. We further validate cross-sensor generalization by transferring trained scenes to a different radar configuration. This opens a path toward physics-grounded radar sensing pipelines, including calibration, radar-aware 3D reconstruction, and joint multi-modal optimization.

\textbf{Limitations and future work.} mmIR is bottlenecked by mesh quality: the LiDAR-to-mesh reconstruction pipeline sets a ceiling on rendered ADC realism, and the inverse rendering task is sensitive to LiDAR-radar alignment for both cascade and single-chip configurations. Our prototype further assumes static scenes and surface-only interactions (no transmission, refraction, or volumetric scattering), models diffraction at single-bounce interactions only, and does not account for Doppler phase from ego-motion. Although our pipeline supports differentiable vertex positions and pose, we found that at 77~GHz ($\lambda \approx 3.9$~mm), sub-millimeter vertex shifts induce large phase changes, making geometry and pose gradients unstable under Monte Carlo estimation with current sample budgets. Shape-from-radar, where geometry is directly recovered from radar measurements, remains a compelling goal but will likely require dense multi-view supervision and regularization strategies beyond what single-frame RA fitting provides. Natural extensions include higher-fidelity meshing, robust alignment, temporally coherent Doppler modeling, multi-bounce diffraction, and richer propagation effects. We discuss these limitations in detail in the supplement (Sec.~\ref{sec:limitations}). Scaling to jointly optimize radar with other sensors and learning priors over radar-consistent materials are promising directions for robust autonomous perception.
